\chapter{SARS-CoV-2 Testing}

\section*{What Is This Test?}
SARS-CoV-2 tests detect current infection or evidence of a previous immune response. Nucleic acid amplification tests, including PCR-based assays, detect viral genetic material, while antigen tests detect viral proteins. Antibody tests detect immune responses and are not the primary test for diagnosing a current infection.

\section*{Alternative Names}
COVID-19 PCR; SARS-CoV-2 NAAT; COVID Antigen Test; Coronavirus Test; COVID Serology.

\section*{Why This Test Is Ordered}
Testing is used for compatible symptoms, exposure assessment, infection-control decisions, high-risk clinical settings, and selected pre-travel or institutional requirements. A negative result does not always exclude early or low-level infection.

\section*{How This Test Is Performed}
The specimen is commonly collected from the nose or another approved respiratory site using a swab. The laboratory performs a molecular or antigen assay according to the specimen and platform. Antibody testing uses blood and is interpreted separately from tests for current infection.

\section*{Preparation, Risks, and Limitations}
No fasting is required. Swabbing may cause brief discomfort or irritation. Test performance depends on timing, specimen quality, assay sensitivity, and the current circulating variants. Follow current local public-health and clinical guidance for repeat testing and isolation decisions.

\section*{Understanding Results}
A positive molecular or antigen result generally indicates current infection, although false positives can occur. A negative result lowers the likelihood of infection but may require repeat testing when symptoms or exposure risk remain high. A positive antibody result indicates immune response but does not reliably establish current infection, infectiousness, or complete protection.

\section*{References}
World Health Organization and Centers for Disease Control and Prevention guidance on SARS-CoV-2 testing.

\clearpage
\section*{Regulatory Status and Authorization Date}
Regulatory status applies to the specific assay, instrument, manufacturer, and intended use rather than to the test name alone. No single approval date can be assigned to this test category. For a commercial assay, verify the current product in the relevant regulator's database: the U.S. Food and Drug Administration (FDA) clearance, approval, or authorization; India's Central Drugs Standard Control Organisation (CDSCO) licence or permission; the European Union In Vitro Diagnostic Regulation (EU IVDR) CE certificate issued through the applicable conformity-assessment route; or the United Kingdom Medicines and Healthcare products Regulatory Agency (MHRA) registration and UKCA/CE status. Laboratory-developed tests may not have an individual FDA, CDSCO, EU IVDR, or MHRA approval date.
\clearpage
